\subsection{Fibres over completed residue fields}
\label{subsec:completed-residue-fibres}

Let $A\to B$ be a homomorphism of arbitrary commutative rings, represented
as an $A$-algebra structure on $B$, and let $x\in\mathcal S(A)$. Write
$H(x)$ for its genuine completed residue field, $C=H(x)\otimes_A B$, and
$v_{H(x)}(c)=\lVert c\rVert$ for its prescribed nonarchimedean norm.
The constructions below require no original coefficient field, trivial norm,
integrality, finite generation, or completeness of $A$ or $B$.

\begin{lemma}[Tensor evaluation at a fibre point]
\label{spectrum:fibre-tensor-evaluation}
\lean{EnvelopeProperty.Berkovich.Spectrum.fibreHom,
EnvelopeProperty.Berkovich.Spectrum.fibreHom_tmul,
EnvelopeProperty.Berkovich.Spectrum.fibreHom_one_tmul,
EnvelopeProperty.Berkovich.Spectrum.fibreHom_tmul_one,
EnvelopeProperty.Berkovich.Spectrum.algebraMap_tmul_one_eq,
EnvelopeProperty.Berkovich.Spectrum.fibrePoint,
EnvelopeProperty.Berkovich.Spectrum.fibrePoint_val_tmul}
\leanok
\uses{spectrum:completed-residue-field,spectrum:completed-residue-field-pullback}
If $y\in\mathcal S(B)$ restricts to $x$, there is an actual ring homomorphism
\[
 e_y:H(x)\otimes_A B\longrightarrow H(y),\qquad
 c\otimes b\longmapsto H(A\to B)(c)\,\iota_y(b).
\]
Pullback of the norm of $H(y)$ gives a real seminorm on $C$ restricting
exactly to the norm of $H(x)$. On elementary tensors its value is
$\lVert c\rVert\,|b|_y$, and restriction along $b\mapsto1\otimes b$
recovers $y$.
\end{lemma}
\begin{proof}
The completed residue-field map preserves norms and commutes with the
maps from $A$. Thus it and $B\to H(y)$ are compatible $A$-algebra maps,
and the tensor universal property constructs $e_y$.
Pullback of the genuine norm valuation is a multiplicative ultrametric
seminorm. Norm preservation on $H(x)$ proves the prescribed restriction;
multiplicativity and norm preservation on $B$ give the tensor formula and
the asserted recovery of $y$.
\end{proof}

\begin{theorem}[Completed-residue-field fibre equivalence]
\label{spectrum:fibre-equivalence}
\lean{EnvelopeProperty.Berkovich.Spectrum.Fibre,
EnvelopeProperty.Berkovich.Spectrum.fibreRestrict,
EnvelopeProperty.Berkovich.Spectrum.fibreRestrict_val,
EnvelopeProperty.Berkovich.Spectrum.continuous_fibreRestrict,
EnvelopeProperty.Berkovich.Spectrum.fibreRestrict_fibrePoint,
EnvelopeProperty.Berkovich.Spectrum.fibreHom_restrict_comp_toH,
EnvelopeProperty.Berkovich.Spectrum.fibrePoint_fibreRestrict,
EnvelopeProperty.Berkovich.Spectrum.fibreEquiv}
\leanok
\uses{spectrum:fibre-tensor-evaluation,spectrum:completed-residue-field-pullback}
Restriction along $B\to C$ gives a continuous map
\[
 \mathcal S_{v_{H(x)}}(C)\longrightarrow
 \{y\in\mathcal S(B):y|_A=x\}.
\]
This map is a bijection. Its inverse assigns to $y$ the norm pulled back
along $e_y$. This assertion is an equivalence of point sets with a continuous
restriction map; inverse continuity is a separate assertion.
\end{theorem}
\begin{proof}
For a prescribed-norm point $z$ of $C$, let $y$ be its restriction to $B$.
The tensor balancing relation identifies $1\otimes a$ with
$\iota_x(a)\otimes1$, and the prescribed norm therefore gives
$|a|_y=|a|_x$. Thus restriction lands in the claimed fibre, continuously
by continuous seminorm pullback.
The preceding lemma proves restriction after tensor evaluation is the identity.
For the other identity, the genuine completed-field map $H(y)\to H(z)$
is an isometry. The two ring maps from $H(x)$ to $H(z)$, respectively
through $H(y)$ and through $C\to H(z)$, agree on $\kappa(x)$:
the residue-field universal property reduces their agreement to elements
of $A$, where it is the tensor balancing relation. Both maps are continuous
and $\kappa(x)$ is dense in $H(x)$, so they agree everywhere.
Consequently the two maps from $C$ to $H(z)$ agree on every elementary
tensor and hence on every tensor element by additivity. Norm preservation
of $H(y)\to H(z)$ and $C\to H(z)$ now proves equality of the entire
seminorms. This proves the second inverse identity.
No existence of a point above an arbitrary point under a separate base
extension is asserted here.
\end{proof}

\begin{lemma}[Limits of coordinate-bounded ultrafilters]
\label{spectrum:bounded-ultrafilter-limit}
\lean{EnvelopeProperty.Berkovich.RelativeSpectrum.exists_ultrafilter_limit_of_eventually_bounded}
\leanok
\uses{spectrum:topology,spectrum:closed-evaluation-embedding,spectrum:relative}
In any prescribed-base real seminorm spectrum, suppose an ultrafilter
has a finite upper bound eventually on each evaluation coordinate,
with bounds and eventual sets allowed to depend on that coordinate.
Then the ultrafilter converges to a point of the same prescribed-base spectrum.
\end{lemma}
\begin{proof}
For each ring element $a$, map the ultrafilter under evaluation at $a$.
Its eventual bound places it in a compact interval $[0,M_a]$, so this
mapped ultrafilter converges. Choose these limits for all coordinates.
The product topology makes the evaluation-function ultrafilter converge
to their product. The seminorm spectrum is closed in the evaluation
product, and the prescribed-base locus is closed in it. Their closed
embedding therefore contains the limit and reflects convergence,
giving a limit in the prescribed-base spectrum itself.
No single eventual set bounding all coordinates at once is required.
\end{proof}

\begin{theorem}[Completed-residue-field fibre homeomorphism]
\label{spectrum:fibre-homeomorphism}
\lean{EnvelopeProperty.Berkovich.Spectrum.fibreRestrict_val_tmul,
EnvelopeProperty.Berkovich.Spectrum.isProperMap_fibreRestrict,
EnvelopeProperty.Berkovich.Spectrum.isClosedMap_fibreRestrict,
EnvelopeProperty.Berkovich.Spectrum.fibreHomeomorph,
EnvelopeProperty.Berkovich.Spectrum.fibreHomeomorph_apply,
EnvelopeProperty.Berkovich.Spectrum.fibreHomeomorph_symm_apply,
EnvelopeProperty.Berkovich.Spectrum.continuous_fibrePoint}
\leanok
\uses{spectrum:bounded-ultrafilter-limit,spectrum:fibre-equivalence,
spectrum:base-change-restriction}
The genuine equivalence of the prescribed-norm tensor spectrum with
the fibre is a homeomorphism for arbitrary commutative $A$-algebras $B$.
No finite-generation or local-compactness hypothesis is required.
\end{theorem}
\begin{proof}
Consider an ultrafilter of tensor points whose restrictions converge
to a point of the original fibre. Every evaluation $|b|$ of those
restrictions is eventually bounded by continuity. The generic prescribed-base
tensor formula gives $|c\otimes b|=\lVert c\rVert|b|$.
Every tensor is a finite sum of elementary tensors, so the ultrametric
inequality supplies an eventual bound for each tensor evaluation.
The preceding lemma gives a convergent tensor point. Continuous restriction
and uniqueness of limits identify its image with the original fibre limit.
This is the ultrafilter lifting criterion for a proper map, so restriction
is closed. It is already a continuous bijection by the fibre equivalence.
Its inverse is therefore continuous, yielding the stated homeomorphism.
\end{proof}
